% GENERATED FILE. Edit canonical lessons, then run npm run generate:books.
% canonical-source-hash: fnv1a64:9b3939099a07c292
% canonical-lessons: KA-C167-uduru, KA-C167-sudharisu, KA-C167-bembalisu, KA-C167-sidi, KA-C167-gorake-hode

\chapter{To drop off, To improve, To support, To burst, To snore}
\label{ch:ka-to-drop-off-to-improve-to-support-to-burst-to-snore}
\hlchaptermodality{167}

\begin{chapteropening}
\textbf{By the end of this chapter:} \emph{I can say to drop off, to improve, to support, to burst and to snore.}

Complete the last lesson of chapter 167: I can say to drop off, to improve, to support, to burst and to snore.
\end{chapteropening}

\section[\texorpdfstring{\kn{ಉದುರು} (uduru)}{ಉದುರು (uduru)}]{\kn{ಉದುರು} (uduru) — to drop off, to shed (leaves, hair)}

\label{lesson:KA-C167-uduru}



\begin{quote}
Before the new one: say the Kannada for to leak, then the Kannada for to float.
\end{quote}

\subsection*{You'll want to know: \kn{ಉದುರು}}
\textbf{\kn{ಉದುರು}} — \emph{uduru} — \textquotedblleft{}to drop off, to shed (leaves, hair)\textquotedblright{}.

\begin{grammarlens}[title={What you've built}]
One more word for this chapter.
\end{grammarlens}

\subsection*{Your turn}
\begin{itemize}
\raggedright
  \item \emph{Say it:} \emph{uduru}
  \item \emph{Say it:} \emph{uduru}, once more
  \item \emph{Say it:} say \emph{uduru}
  \item \emph{Recall:} say the Kannada for to leak, then the Kannada for to float, then say \emph{uduru} again
\end{itemize}

\begin{tcolorbox}[breakable,colback=teal!4,colframe=teal!35!black,title={Before you move on}]
What does \kn{ಉದುರು} mean? (\textquotedblleft{}To drop off, to shed (leaves, hair)\textquotedblright{}.) Say it once more.
\end{tcolorbox}
\section[\texorpdfstring{\kn{ಸುಧಾರಿಸು} (sudhārisu)}{ಸುಧಾರಿಸು (sudhārisu)}]{\kn{ಸುಧಾರಿಸು} (sudhārisu) — to improve}

\label{lesson:KA-C167-sudharisu}



\begin{quote}
Before the new one: say the Kannada for to float, then the Kannada for to drop off.
\end{quote}

\subsection*{You'll want to know: \kn{ಸುಧಾರಿಸು}}
\textbf{\kn{ಸುಧಾರಿಸು}} — \emph{sudhārisu} — \textquotedblleft{}to improve\textquotedblright{}.

\begin{grammarlens}[title={What you've built}]
One more word for this chapter.
\end{grammarlens}

\subsection*{Your turn}
\begin{itemize}
\raggedright
  \item \emph{Say it:} \emph{sudhārisu}
  \item \emph{Say it:} \emph{sudhārisu}, once more
  \item \emph{Say it:} say \emph{sudhārisu}
  \item \emph{Recall:} say the Kannada for to float, then the Kannada for to drop off, then say \emph{sudhārisu} again
\end{itemize}

\begin{tcolorbox}[breakable,colback=teal!4,colframe=teal!35!black,title={Before you move on}]
What does \kn{ಸುಧಾರಿಸು} mean? (\textquotedblleft{}To improve\textquotedblright{}.) Say it once more.
\end{tcolorbox}
\section[\texorpdfstring{\kn{ಬೆಂಬಲಿಸು} (bembalisu)}{ಬೆಂಬಲಿಸು (bembalisu)}]{\kn{ಬೆಂಬಲಿಸು} (bembalisu) — to support, to back}

\label{lesson:KA-C167-bembalisu}



\begin{quote}
Before the new one: say the Kannada for to drop off, then the Kannada for to improve.
\end{quote}

\subsection*{You'll want to know: \kn{ಬೆಂಬಲಿಸು}}
\textbf{\kn{ಬೆಂಬಲಿಸು}} — \emph{bembalisu} — \textquotedblleft{}to support, to back\textquotedblright{}.

\begin{grammarlens}[title={What you've built}]
One more word for this chapter.
\end{grammarlens}

\subsection*{Your turn}
\begin{itemize}
\raggedright
  \item \emph{Say it:} \emph{bembalisu}
  \item \emph{Say it:} \emph{bembalisu}, once more
  \item \emph{Say it:} say \emph{bembalisu}
  \item \emph{Recall:} say the Kannada for to drop off, then the Kannada for to improve, then say \emph{bembalisu} again
\end{itemize}

\begin{tcolorbox}[breakable,colback=teal!4,colframe=teal!35!black,title={Before you move on}]
What does \kn{ಬೆಂಬಲಿಸು} mean? (\textquotedblleft{}To support, to back\textquotedblright{}.) Say it once more.
\end{tcolorbox}
\section[\texorpdfstring{\kn{ಸಿಡಿ} (siḍi)}{ಸಿಡಿ (siḍi)}]{\kn{ಸಿಡಿ} (siḍi) — to burst, to explode}

\label{lesson:KA-C167-sidi}



\begin{quote}
Before the new one: say the Kannada for to improve, then the Kannada for to support.
\end{quote}

\subsection*{You'll want to know: \kn{ಸಿಡಿ}}
\textbf{\kn{ಸಿಡಿ}} — \emph{siḍi} — \textquotedblleft{}to burst, to explode\textquotedblright{}.

\begin{grammarlens}[title={What you've built}]
One more word for this chapter.
\end{grammarlens}

\subsection*{Your turn}
\begin{itemize}
\raggedright
  \item \emph{Say it:} \emph{siḍi}
  \item \emph{Say it:} \emph{siḍi}, once more
  \item \emph{Say it:} say \emph{siḍi}
  \item \emph{Recall:} say the Kannada for to improve, then the Kannada for to support, then say \emph{siḍi} again
\end{itemize}

\begin{tcolorbox}[breakable,colback=teal!4,colframe=teal!35!black,title={Before you move on}]
What does \kn{ಸಿಡಿ} mean? (\textquotedblleft{}To burst, to explode\textquotedblright{}.) Say it once more.
\end{tcolorbox}
\section[\texorpdfstring{\kn{ಗೊರಕೆ} \kn{ಹೊಡೆ} (gorake hoḍe)}{ಗೊರಕೆ ಹೊಡೆ (gorake hoḍe)}]{\kn{ಗೊರಕೆ} \kn{ಹೊಡೆ} (gorake hoḍe) — to snore}

\label{lesson:KA-C167-gorake-hode}



\begin{quote}
Before the new one: say the Kannada for to support, then the Kannada for to burst.
\end{quote}

\subsection*{You'll want to know: \kn{ಗೊರಕೆ} \kn{ಹೊಡೆ}}
\textbf{\kn{ಗೊರಕೆ} \kn{ಹೊಡೆ}} — \emph{gorake hoḍe} — \textquotedblleft{}to snore\textquotedblright{}.

\begin{grammarlens}[title={What you've built}]
One more word for this chapter.
\end{grammarlens}

\subsection*{Your turn}
\begin{itemize}
\raggedright
  \item \emph{Say it:} \emph{gorake hoḍe}
  \item \emph{Say it:} \emph{gorake hoḍe}, once more
  \item \emph{Say it:} say \emph{gorake hoḍe}
  \item \emph{Recall:} say the Kannada for to support, then the Kannada for to burst, then say \emph{gorake hoḍe} again
\end{itemize}

\begin{tcolorbox}[breakable,colback=teal!4,colframe=teal!35!black,title={Before you move on}]
What does \kn{ಗೊರಕೆ} \kn{ಹೊಡೆ} mean? (\textquotedblleft{}To snore\textquotedblright{}.) Say it once more.
\end{tcolorbox}
